% !PRINTED 379
\weberpage{379}
%<<CONT
Let $\varepsilon$ run through all the divisors of $m$ (1 and $m$
included) and $e^2$ through all the quadratic divisors of
$m$; then

\weq{10}{\overset{e^2}{\Sigma}\ \psi\!\left(D,\frac{m}{e^2}\right) = \overset{\varepsilon}{\Sigma}\ (D,\varepsilon).}

If this relation is correct for two numbers $m'$, $m''$ that have no
common divisor, then it follows for $m=m'm''$. For if
$e'$, $\varepsilon'$ and $e''$, $\varepsilon''$ have the same meaning for $m'$ and $m''$ as
$e$, $\varepsilon$ have for $m$, then $\varepsilon=\varepsilon'\varepsilon''$ are the divisors and $e^2=e'^2e''^2$ the
quadratic divisors of $m=m'm''$. But by (6) and
\S\,85, (8) we have:

\[ \overset{e'^2}{\Sigma}\ \overset{e''^2}{\Sigma}\ \psi\!\left(D,\frac{m'}{e'^2}\right)\psi\!\left(D,\frac{m''}{e''^2}\right) = \overset{e^2}{\Sigma}\ \psi\!\left(D,\frac{m}{e^2}\right), \]

\[ \overset{\varepsilon'}{\Sigma}\ \overset{\varepsilon''}{\Sigma}\ (D,\varepsilon')(D,\varepsilon'') = \Sigma(D,\varepsilon). \]

Accordingly we need only prove the relation (10) under the
assumption that $m=q^k$ is a prime power.

But under this assumption we have:

\[
\begin{aligned}
\frac{n}{e^2}&=q^{k-2s}, & 0&\leqq s \leqq \frac{k}{2},\\
\varepsilon&=q^s, & s&=0,1,2,\dots,k.
\end{aligned}
\] % [sic] German source prints n instead of m; the second relation uses a visibly different sort.

If now

1. $(D,q)=0$, then $\overset{s}{\Sigma}\ \psi(D,q^{k-2s}) = 1$,

2. $(D,q)=-1$, \text{\glqq}\ \text{\glqq}\ $\overset{s}{\Sigma}\ \psi(D,q^{k-2s}) = 1$ ($k$ even),
\hspace{4em}$= 0$ ($k$ odd),

3. $(D,q)=1$, \text{\glqq}\ \text{\glqq}\ $\overset{s}{\Sigma}\ \psi(D,q^{k-2s}) = k+1$,

and in the three cases just as large a value is obtained for the sum

\[ \underset{0,\,k}{\overset{s}{\Sigma}}\ (D,q^s); \]

for in case 1. only the term of this sum corresponding to $s=0$
has the value 1, the others vanish; in case 2.
the terms corresponding to the even $s$ have the value $+1$,
the others the value $-1$, and in case 3. all
$k+1$-terms have the value $+1$.

The relation (10) is thus proved in general.
